\section{Estimación por intervalos de confianza}

\begin{ejemplo}
    Sea $X_1, \ldots, X_n$ una muestra aleatoria simple de una variable aleatoria $X \rightsquigarrow \cc{N}(\mu, \sigma^2_0)$, con $\mu \in \mathbb{R}$. Construir un intervalo de confianza para $\mu$ al nivel de confianza $1 - \alpha$ de la forma $\left]\ol{X} - C_1, \ol{X} + C_2\right[$.\\

    Buscamos obtener $C_1$ y $C_2$ tales que:
    \begin{align*}
        P_\mu[\ol{X} - C_1 \leq \mu \leq \ol{X} + C_2] \geq 1 - \alpha
    \end{align*}

    La única variable aleatoria que aparece en la desigualdad es $\ol{X}$, por lo que podemos reescribirla como:
    \begin{align*}
        P_\mu[\mu - C_2 \leq \ol{X} \leq \mu + C_1] \geq 1 - \alpha
    \end{align*}

    Tipificando $\ol{X}$ sabiendo que $\ol{X} \rightsquigarrow \cc{N}\left(\mu, \nicefrac{\sigma^2_0}{n}\right)$, obtenemos:
    \begin{align*}
        P_\mu\left[\frac{\mu - C_2 - \mu}{\nicefrac{\sigma_0}{\sqrt{n}}} \leq \frac{\ol{X} - \mu}{\nicefrac{\sigma_0}{\sqrt{n}}} \leq \frac{\mu + C_1 - \mu}{\nicefrac{\sigma_0}{\sqrt{n}}}\right] \geq 1 - \alpha
    \end{align*}

    Notando por $Z$ a la variable aleatoria tipificada de $\ol{X}$, $Z \rightsquigarrow \cc{N}(0, 1)$, podemos reescribir la desigualdad como:
    \begin{align*}
        P_\mu\left[\frac{- C_2}{\nicefrac{\sigma_0}{\sqrt{n}}} \leq Z \leq \frac{C_1}{\nicefrac{\sigma_0}{\sqrt{n}}}\right] \geq 1 - \alpha
    \end{align*}

    Aunque habría muchas formas de elegir $\frac{C_1}{\nicefrac{\sigma_0}{\sqrt{n}}}$ y $\frac{- C_2}{\nicefrac{\sigma_0}{\sqrt{n}}}$, una forma natural de hacerlo es que sean simétricos respecto a $0$, pues será el intervalo de confianza más pequeño posible. Por lo tanto, podemos imponer que:
    \begin{align*}
        \frac{C_1}{\nicefrac{\sigma_0}{\sqrt{n}}} = \frac{- C_2}{\nicefrac{\sigma_0}{\sqrt{n}}} \Longrightarrow C_1 = C_2 = C
    \end{align*}

    En ese caso, tenemos que:
    \begin{align*}
        \dfrac{C}{\nicefrac{\sigma_0}{\sqrt{n}}} = z_{\nicefrac{\alpha}{2}}
    \end{align*}
    donde $P[Z\geq z_{\nicefrac{\alpha}{2}}] = \nicefrac{\alpha}{2}$. Por lo tanto, obtenemos que:
    \begin{align*}
        C = z_{\nicefrac{\alpha}{2}} \dfrac{\sigma_0}{\sqrt{n}}
    \end{align*}

    Por lo tanto, el intervalo de confianza para $\mu$ al nivel de confianza $1 - \alpha$ es:
    \begin{align*}
        \left]\ol{X} - z_{\nicefrac{\alpha}{2}} \dfrac{\sigma_0}{\sqrt{n}}, \ol{X} + z_{\nicefrac{\alpha}{2}} \dfrac{\sigma_0}{\sqrt{n}}\right[
    \end{align*}

    En el caso típico de $\alpha = 0.05$, tenemos que $z_{\nicefrac{\alpha}{2}} = z_{0.025} = 1.96$, por lo que el intervalo de confianza para $\mu$ al nivel de confianza $95\%$ es:
    \begin{align*}
        \left]\ol{X} - 1.96 \dfrac{\sigma_0}{\sqrt{n}}, \ol{X} + 1.96 \dfrac{\sigma_0}{\sqrt{n}}\right[
    \end{align*}
        
\end{ejemplo}



\begin{ejemplo}
    Obtener, utilizando la desigualdad de Chebychev, un intervalo de confianza para la media poblacional, $\mu$, basado en una muestra de tamaño $n$ de una variable aleatoria cuya varianza, $\sigma^2_0$, es conocida.\\

    Sea $X_1, \ldots, X_n$ una muestra aleatoria simple de una variable aleatoria $X \rightsquigarrow \cc{N}(\mu, \sigma^2_0)$, con $\mu \in \mathbb{R}$. Sabemos que $\ol{X} \rightsquigarrow \cc{N}\left(\mu, \nicefrac{\sigma^2_0}{n}\right)$, por lo que:
    \begin{align*}
        E[\ol{X}] = \mu, \quad Var[\ol{X}] = \frac{\sigma^2_0}{n}
    \end{align*}

    Por tanto, por la desigualdad de Chebychev, para cualquier $k > 0$:
    \begin{align*}
        P_\mu[|\ol{X} - \mu| \geq k]\geq 1-  \frac{Var[\ol{X}]}{k^2} = 1-\frac{\sigma^2_0}{n k^2}
    \end{align*}

    Reescribiendo la desigualdad, obtenemos:
    \begin{align*}
        P_\mu[k-\mu \leq \ol{X} \leq k+\mu] \geq 1 - \frac{\sigma^2_0}{n k^2}
    \end{align*}

    Como en un intervalo de confianza queremos calcular la probabilidad de que el intervalo contenga a $\mu$, podemos reescribir la desigualdad como:
    \begin{align*}
        P_\mu[\ol{X} - k \leq \mu \leq \ol{X} + k] \geq 1 - \frac{\sigma^2_0}{n k^2}
    \end{align*}

    Por lo tanto, para obtener un intervalo de confianza al nivel de confianza $1 - \alpha$, debemos imponer que:
    \begin{align*}
        1 - \frac{\sigma^2_0}{n k^2} = 1 - \alpha \Longrightarrow k = \sqrt{\frac{\sigma^2_0}{n \alpha}}
    \end{align*}

    Por lo tanto, el intervalo de confianza para $\mu$ al nivel de confianza $1 - \alpha$ es:
    \begin{align*}
        \left]\ol{X} - \sqrt{\frac{\sigma^2_0}{n \alpha}}, \ol{X} + \sqrt{\frac{\sigma^2_0}{n \alpha}}\right[
    \end{align*}
\end{ejemplo}